% Zumdahl Chapter 16 test -- 20 MCQ + 2 short FRQ. 50 min.
\documentclass{apchem-exam}

\apcunitword{Chapter}
\apcunit{16}
\apcunittitle{Solubility Equilibria}
\apcdoctitle{Chapter Test}
\apcsubtitle{Zumdahl \S16.1 \textbullet\ 50 minutes \textbullet\ $K_{sp}$ values provided with Section II}

\begin{document}
\apctitleblock

\begin{apcnote}[Scope]
Covers \S16.1 only --- CED topics 7.11, 7.12 and 8.11. Nothing from
\S16.2 (precipitation, qualitative analysis) or \S16.3 (complex ions)
appears; both are off the framework, and the $Q$-versus-$K_{sp}$ material in
Block 4 was marked enrichment.

Per the CED exclusion on topic 8.11, pH and solubility is assessed
\textbf{qualitatively only}. You will be asked to predict a direction and
justify it, never to compute a solubility from a pH.
\end{apcnote}

\examsection{I}{Multiple Choice}{20 questions \textbullet\ 25 minutes \textbullet\ 1 point each}{\nocalculator}

% ---- Ksp expressions and meaning -------------------------------------------
\begin{mcq}
  The $K_{sp}$ expression for \ce{Ag2CrO4} is
  \choices
    \choice{$[\ce{Ag+}][\ce{CrO4^2-}]$}
    \correct{$[\ce{Ag+}]^2[\ce{CrO4^2-}]$}
    \choice{$[\ce{Ag+}]^2[\ce{CrO4^2-}]^2$}
    \choice{$[\ce{Ag2CrO4}]/[\ce{Ag+}]^2[\ce{CrO4^2-}]$}
\end{mcq}
\why{Each ion is raised to its coefficient; the solid is omitted.}

\begin{mcq}
  A solid is omitted from a $K_{sp}$ expression because
  \choices
    \correct{it is a pure solid and has no variable concentration}
    \choice{its concentration is zero}
    \choice{it does not participate in the reaction}
    \choice{$K_{sp}$ applies only to ions}
\end{mcq}
\why{The same rule that omits pure solids and liquids from any $K$.}

\begin{mcq}
  Adding more undissolved solid \ce{AgCl} to a saturated solution of
  \ce{AgCl} causes $[\ce{Ag+}]$ to
  \choices
    \correct{stay the same}
    \choice{increase}
    \choice{decrease}
    \choice{increase, then decrease}
\end{mcq}
\why{Extra surface area speeds dissolving and re-depositing equally, so the
equilibrium position does not move.}

\begin{mcq}
  Which statement is correct?
  \choices
    \correct{$K_{sp}$ is a constant; solubility is a position that varies
      with the solution}
    \choice{$K_{sp}$ and solubility are two names for the same quantity}
    \choice{solubility is constant; $K_{sp}$ varies with the solution}
    \choice{both vary with the amount of excess solid}
\end{mcq}
\why{This distinction governs every common-ion question.}

% ---- Ksp from solubility and back ------------------------------------------
\begin{mcq}
  For a salt \ce{MX} with molar solubility $s$, $K_{sp}$ equals
  \choices
    \correct{$s^2$}
    \choice{$2s$}
    \choice{$4s^3$}
    \choice{$s$}
\end{mcq}
\why{Both ions are present at $s$.}

\begin{mcq}
  For a salt \ce{MX2} with molar solubility $s$, $K_{sp}$ equals
  \choices
    \choice{$s^3$}
    \choice{$2s^3$}
    \correct{$4s^3$}
    \choice{$27s^4$}
\end{mcq}
\why{$(s)(2s)^2$ --- the 2 is both a multiplier and inside the square.}
\distractor{A}{forgot to double the anion}

\begin{mcq}
  A salt \ce{MX} has $K_{sp} = 4.0\times10^{-10}$. Its molar solubility is
  \choices
    \choice{$4.0\times10^{-10}$}
    \choice{$2.0\times10^{-10}$}
    \correct{$2.0\times10^{-5}$}
    \choice{$1.6\times10^{-19}$}
\end{mcq}
\why{$s = \sqrt{K_{sp}}$.}

\begin{mcq}
  A student finds $K_{sp}$ for \ce{PbI2} using $K_{sp} = s^3$. Their answer
  is
  \choices
    \correct{too small by a factor of 4}
    \choice{too large by a factor of 4}
    \choice{too small by a factor of 2}
    \choice{correct}
\end{mcq}
\why{The correct expression is $4s^3$.}

\begin{mcq}
  \ce{Ag3PO4} dissolves to give three \ce{Ag+} and one \ce{PO4^3-}. In
  terms of $s$, $K_{sp}$ equals
  \choices
    \choice{$3s^2$}
    \choice{$4s^3$}
    \correct{$27s^4$}
    \choice{$s^4$}
\end{mcq}
\why{$(3s)^3(s) = 27s^4$.}

% ---- relative solubilities -------------------------------------------------
\begin{mcq}
  $K_{sp}$ values may be compared directly to rank solubilities only when
  the salts
  \choices
    \correct{produce the same total number of ions}
    \choice{contain the same cation}
    \choice{have $K_{sp}$ values within one order of magnitude}
    \choice{are all insoluble by the Unit 4 rules}
\end{mcq}
\why{Otherwise $s$ enters $K_{sp}$ raised to different powers.}

\begin{mcq}
  \ce{BaSO4} ($K_{sp} = 1.5\times10^{-9}$), \ce{PbSO4}
  ($1.3\times10^{-8}$) and \ce{SrSO4} ($3.2\times10^{-7}$). The most
  soluble is
  \choices
    \choice{\ce{BaSO4}}
    \choice{\ce{PbSO4}}
    \correct{\ce{SrSO4}}
    \choice{cannot be determined without calculation}
\end{mcq}
\why{All three are \ce{MX} salts, so the largest $K_{sp}$ is the most
soluble.}

\begin{mcq}
  \ce{CuS} has a much larger $K_{sp}$ than \ce{Bi2S3}, yet \ce{Bi2S3} is
  far more soluble. This is possible because
  \choices
    \correct{\ce{Bi2S3} releases five ions, so $K_{sp}$ involves $s$ to the
      fifth power}
    \choice{$K_{sp}$ values for sulfides are unreliable}
    \choice{\ce{Bi2S3} reacts with water}
    \choice{\ce{CuS} has a larger molar mass}
\end{mcq}
\why{$K_{sp} = 108s^5$ for \ce{M2X3}; a fifth power crushes the number.}

\begin{mcq}
  According to the CED, salts the Unit 4 rules call ``soluble'' correspond
  to $K_{sp}$ values that are
  \choices
    \choice{negative}
    \choice{exactly 1}
    \correct{greater than 1}
    \choice{smaller than $10^{-10}$}
\end{mcq}
\why{The solubility rules are the qualitative form of the same idea.}

% ---- common ion ------------------------------------------------------------
\begin{mcq}
  Adding \ce{NaCl} to a saturated solution of \ce{AgCl} causes the
  solubility of \ce{AgCl} to
  \choices
    \correct{decrease}
    \choice{increase}
    \choice{remain the same}
    \choice{decrease only if the solution is heated}
\end{mcq}
\why{Adding a product shifts the dissolution equilibrium left.}

\begin{mcq}
  In that same experiment, $K_{sp}$ for \ce{AgCl}
  \choices
    \choice{decreases}
    \choice{increases}
    \correct{does not change}
    \choice{becomes undefined}
\end{mcq}
\why{Only temperature changes an equilibrium constant.}

\begin{mcq}
  Which addition would \emph{not} change the solubility of \ce{BaSO4}?
  \choices
    \choice{\ce{Na2SO4}}
    \choice{\ce{BaCl2}}
    \correct{\ce{KNO3}}
    \choice{\ce{Ba(NO3)2}}
\end{mcq}
\why{Neither \ce{K+} nor \ce{NO3-} appears in the equilibrium.}

\begin{mcq}
  For \ce{CaF2}, adding fluoride suppresses the solubility far more
  strongly than adding an equal concentration of calcium, because
  \choices
    \correct{fluoride is squared in the $K_{sp}$ expression}
    \choice{fluoride ions are smaller}
    \choice{calcium ions are more highly charged}
    \choice{\ce{NaF} is more soluble than \ce{Ca(NO3)2}}
\end{mcq}
\why{The exponent controls how sensitive the product is to that ion.}

% ---- pH and solubility -----------------------------------------------------
\begin{mcq}
  A salt's solubility increases in acidic solution when its anion is
  \choices
    \correct{the conjugate base of a weak acid}
    \choice{the conjugate base of a strong acid}
    \choice{a halide}
    \choice{singly charged}
\end{mcq}
\why{Only then does added \ce{H+} remove the anion from solution.}

\begin{mcq}
  Which salt is \emph{more} soluble in \SI{1}{\Molar} \ce{HNO3} than in
  pure water?
  \choices
    \choice{\ce{AgCl}}
    \choice{\ce{AgBr}}
    \correct{\ce{CaCO3}}
    \choice{\ce{PbI2}}
\end{mcq}
\why{\ce{CO3^2-} is the conjugate base of the weak acid \ce{H2CO3}; the
halides come from strong acids.}
\distractor{A}{\ce{Cl-} is a negligible base, so acid has no effect}

\begin{mcq}
  Raising the pH of a saturated \ce{Mg(OH)2} solution causes its solubility
  to
  \choices
    \choice{increase}
    \correct{decrease}
    \choice{stay the same}
    \choice{first increase, then decrease}
\end{mcq}
\why{Higher pH means more \ce{OH-}, which is a common ion here.}

\examsection{II}{Free Response}{2 questions \textbullet\ 25 minutes}{\calculatorok}

\begin{apcnote}[Data]
$K_{sp}$ at \SI{25}{\celsius}: \ce{Mg(OH)2} $8.9\times10^{-12}$
\textbullet\ \ce{AgCl} $1.6\times10^{-10}$ \textbullet\ \ce{BaSO4}
$1.5\times10^{-9}$ \textbullet\ \ce{CaCO3} $8.7\times10^{-9}$.
\end{apcnote}

% ---- FRQ 1 -----------------------------------------------------------------
\frq{short}{4}{Zumdahl \S16.1 \textbullet\ $K_{sp}$ and molar solubility}

Magnesium hydroxide, \ce{Mg(OH)2}, is a sparingly soluble solid.

\begin{parts}
  \item Write the dissolution equation and the $K_{sp}$ expression.
        \answerlines[2]{\ce{Mg(OH)2(s) <=> Mg^2+(aq) + 2OH^-(aq)};
        $K_{sp} = [\ce{Mg^2+}][\ce{OH-}]^2$}
  \item Calculate the molar solubility in pure water. \workspace[2.6cm]
        \keyonly{\begin{apcanswer}$[\ce{Mg^2+}] = s$, $[\ce{OH-}] = 2s$;
        $K_{sp} = 4s^3 = 8.9\times10^{-12}$;
        $s^3 = 2.2\times10^{-12}$;
        $s = \mathbf{1.3\times10^{-4}}$~M\end{apcanswer}}
  \item State what happens to the solubility if \ce{NaOH} is added, and
        what happens to $K_{sp}$. \answerlines[2]{The solubility decreases,
        because \ce{OH-} is a common ion and shifts the equilibrium left.
        $K_{sp}$ is unchanged --- it depends only on temperature.}
\end{parts}

\begin{rubric}
  \rpoint{1}{(a) equation and expression both correct, with the hydroxide
    squared}
  \rpoint{2}{(b) 1 pt recognizing $[\ce{OH-}] = 2s$ and building $4s^3$;
    1 pt $s = 1.3\times10^{-4}$~M --- using $s^3$ caps this part at 1}
  \rpoint{1}{(c) BOTH halves right: solubility falls AND $K_{sp}$ unchanged
    --- saying $K_{sp}$ decreases earns 0 for this point}
\end{rubric}

% ---- FRQ 2 -----------------------------------------------------------------
\frq{short}{4}{Zumdahl \S16.1 \textbullet\ common ion and pH}

A saturated solution of \ce{BaSO4} is prepared.

\begin{parts}
  \item Calculate the molar solubility of \ce{BaSO4} in pure water.
        \workspace[2cm]
        \keyonly{\begin{apcanswer}$K_{sp} = s^2 = 1.5\times10^{-9}$;
        $s = \mathbf{3.9\times10^{-5}}$~M\end{apcanswer}}
  \item Calculate the molar solubility in \SI{0.010}{\Molar} \ce{Na2SO4}.
        \workspace[2.4cm]
        \keyonly{\begin{apcanswer}$[\ce{SO4^2-}] \approx 0.010$;
        $1.5\times10^{-9} = (s)(0.010)$;
        $s = \mathbf{1.5\times10^{-7}}$~M\end{apcanswer}}
  \item \ce{BaSO4} and \ce{CaCO3} are each treated with dilute
        \ce{HNO3}. Predict the effect on each solubility and justify the
        difference. Do not calculate. \answerlines[3]{\ce{CaCO3} becomes
        more soluble: \ce{CO3^2-} is the conjugate base of the weak acid
        \ce{H2CO3}, so \ce{H+} removes it and the equilibrium shifts right.
        \ce{BaSO4} is essentially unaffected, because \ce{SO4^2-} comes
        from the strong acid \ce{H2SO4} and is a negligible base.}
\end{parts}

\begin{rubric}
  \rpoint{1}{(a) $s = 3.9\times10^{-5}$~M}
  \rpoint{1}{(b) $s = 1.5\times10^{-7}$~M from $K_{sp}/[\ce{SO4^2-}]$}
  \rpoint{2}{(c) 1 pt \ce{CaCO3} more soluble with the weak-acid
    justification; 1 pt \ce{BaSO4} unaffected because \ce{SO4^2-} is a
    negligible base --- an answer that only says ``acids dissolve things''
    earns 0}
\end{rubric}

\examtotal

\end{document}
